\documentclass[10pt,a4paper]{article}
\usepackage[utf8]{inputenc}
\usepackage[T1]{fontenc}
\usepackage[english]{babel}
\usepackage{lmodern}
\usepackage[a4paper,margin=1.9cm,top=1.2cm,bottom=1.8cm]{geometry}
\usepackage{graphicx}
\usepackage{booktabs}
\usepackage{float}
\usepackage{xcolor}
\usepackage[hidelinks]{hyperref}
\usepackage{xurl}
\usepackage[font=small,labelfont=bf,skip=3pt]{caption}
\usepackage{parskip}
\setlength{\parskip}{4pt}
\usepackage{titlesec}
\titlespacing*{\section}{0pt}{8pt}{3pt}
\titlespacing*{\subsection}{0pt}{6pt}{2pt}
\graphicspath{{../figures/en/}}
\newcommand{\m}{$-$}

\begin{document}

\begin{center}
{\Large\bfseries Net Position of the Hungarian Power System}\\[3pt]
Author: Veronika Tomažič
\end{center}

\begin{itemize}\setlength\itemsep{2pt}
\item \textbf{a)} Hungary is a permanent net importer of electricity: since 2015 it has imported on average
1000–1400\,MW, about a quarter of its consumption. Almost half of domestic electricity comes from the Paks nuclear
power plant, which runs at a constant output. Gas-fired plants run when electricity is expensive, while solar
capacity grows every year and almost eliminates imports around midday. Hungary buys electricity in the north
(Slovakia, Austria) and sells part of it on to the south.
\item \textbf{b)} For week 47 I forecast an average net position of \textbf{\m1873\,MW} (imports, $\approx$\,\m315\,GWh),
with a range of \textbf{\m2146 to \m1483\,MW} (80\,\%). The forecast was selected from 18 models tested on 202 weeks.
\item \textbf{c)} In my view, the average day-ahead price in week 47 will be \textbf{lower} than in week 45
(290.98\,€/MWh), at about \textbf{260–275\,€/MWh}. The main reason is much stronger wind in Germany, which makes
electricity cheaper across the whole region, while the gas price is unchanged.
\end{itemize}

The analysis is based on public data from 2015 to 15 Nov 2026 at 15-minute resolution: generation by plant type,
consumption, cross-border trading and day-ahead prices (Energy-Charts, which compiles official ENTSO-E data),
weather for five large Hungarian cities and wind in northern Germany (Open-Meteo), and the TTF gas price
(investing.com). All times are in Central European Time (CET).

% =====================================================================
\subsection*{a) Hungary's net position}

The net position is the difference between generation and consumption and equals the difference between
exports and imports. The sign follows the task: \textbf{+ means export, \m\ means import}.

\subsection*{What generates the power}

Hungary has a few main types of power plants that operate very differently (Table~\ref{tab:teh}). The Paks
nuclear power plant produces almost half of domestic electricity and runs at a constant output. Gas-fired plants
run when electricity is expensive, for example in the evening and in winter. When electricity is cheap, they are
switched off and electricity is imported instead, so their output changes the most. Solar plants produce only
during the day and in good weather, but their electricity is almost free. The Mátra lignite plant is being
gradually phased out, and there are few wind and hydro plants in Hungary.

\begin{table}[H]
\centering\footnotesize
\footnotesize\setlength{\tabcolsep}{4pt}
\begin{tabular}{lrrll}
\toprule
type & TWh & share & operation & cost / role in price \\
\midrule
nuclear (Paks) & 15.2 & 45\,\% & constant; outages during maintenance & low \\
gas & 7.4 & 22\,\% & flexible, runs when prices are high & high; often sets the price \\
solar & 6.8 & 20\,\% & daytime only, weather dependent & almost zero; lowers the price at midday \\
lignite & 1.5 & 5\,\% & stable, declining & medium \\
biomass, wind, hydro, other & 2.9 & 8\,\% & small and stable & low \\
\midrule
\textbf{generation / consumption / net} & \multicolumn{4}{l}{\textbf{34.0 / 43.2 / \m9.5 TWh} (imports $\approx$ 22\,\% of consumption)} \\
\bottomrule
\end{tabular}
\caption{Types of power plants (all plants of the same type combined), 2025.}
\label{tab:teh}
\end{table}

\begin{figure}[H]
\centering
\includegraphics[width=0.7\linewidth]{a1b_generation_load_2020_2026.pdf}
\caption{Generation by plant type and consumption, annual averages 2020–2026.}
\label{fig:trend}
\end{figure}

Figure~\ref{fig:trend} shows how this has changed since 2020; earlier years are not shown because solar plants
are missing from the data before then. The gap between the bars and the consumption line is the import. Lignite
declines every year (from 681\,MW in 2015 to 160\,MW in 2026), while solar grows every year (from 181\,MW in 2020 to
946\,MW in 2026). This year solar produces almost as much as gas. Paks was almost constant for 10 years
($\sim$1720\,MW), but this year it is lower (1558\,MW) because it produced less in summer due to the low level of
the Danube. Imports have stayed similar throughout, between 1000 and 1400\,MW (also in 2015–2019), because solar
has replaced the loss of lignite and Paks.


\subsection*{What drives consumption}

Consumption is almost the same over the years, about 4900\,MW on average, but it changes considerably from day
to day. It depends mainly on temperature and on whether it is a working day or a weekend.

When it is warm (above $\sim$15\,°C), temperature has almost no effect on consumption. When it is colder, people
heat, so every degree lower increases consumption by $\sim$95\,MW (Figure~\ref{fig:temp}, right panel:
2022–2026). In summer, when it is very hot (above $\sim$22\,°C), cooling increases consumption again. If the effect
of temperature were estimated with one straight line over the whole year, it would be only $\sim$47\,MW per
degree and the effect of cold weather would be underestimated.

Figure~\ref{fig:temp} shows three periods: all years (2015–2026), years with solar data (2020–2026) and years
with the best data (2022–2026). The pattern is the same in all of them, but the effect of cold weather is
increasing (from $\sim$66 to $\sim$94\,MW per degree), probably because of the growing number of heat pumps.

At weekends consumption is lower because factories and businesses are closed: by $\sim$9\,\% on Saturday and
by $\sim$13\,\% on Sunday.

\begin{figure}[H]
\centering
\includegraphics[width=0.7\linewidth]{a5a_load_temperature.pdf}
\caption{Daily consumption against daily temperature (working days without holidays) for three periods. Straight
line (dashed) against a piecewise line that separates heating, a neutral zone and cooling.}
\label{fig:temp}
\end{figure}

\subsection*{Which borders it trades over}

Hungary is connected to all seven neighbouring countries. Table~\ref{tab:meje} shows the average trade over each
border per year (+ export, \m\ import). Slovakia and Austria are negative in every year: Hungary buys electricity
from them, most of it from Slovakia ($\sim$1000–1250\,MW). Croatia is positive in every year: Hungary sells
electricity there. For Romania, Serbia and Slovenia the direction changes from year to year. Hungary thus buys
electricity in the north, uses part of it at home and sells part of it on to the south.

While reviewing the data I noticed that the figures for individual borders are not trades between two
countries, but are calculated by the algorithm of the common European market. Their sum, the net position, is
therefore the most reliable figure.

\begin{table}[H]
\centering\footnotesize
\begin{tabular}{lrrrrrrrr}
\toprule
 & Slovakia & Austria & Romania & Croatia & Serbia & Slovenia & Ukraine & \textbf{total} \\
\midrule
2024 & \m1029 & \m893 & +378 & +281 & \m38 & \m153 & +185 & \textbf{\m1268} \\
2025 & \m1246 & \m746 & +612 & +350 & \m82 & \m34 & +65 & \textbf{\m1082} \\
2026 & \m1099 & \m495 & \m65 & +335 & \m195 & +102 & +186 & \textbf{\m1233} \\
\bottomrule
\end{tabular}
\caption{Net trade by border, MW (annual averages; + export from HU, \m\ import into HU). 2026: until 15 Nov.}
\label{tab:meje}
\end{table}

\subsection*{How it changes during the day and the year}

\textbf{During the year:} in winter Hungary imports about twice as much electricity as in spring and summer. In
summer, solar plants reduce imports; in winter there is almost no sun and consumption is higher due to heating,
so gas-fired plants run more and the rest is imported.

\textbf{During the day:} since I am forecasting a week in November, I looked at the daily pattern for October and
November in previous years (Figure~\ref{fig:profil}). At midday, when solar plants produce, imports are now
almost zero, whereas in 2020 they were still $\sim$1300\,MW. In the evening, when there is no sun and consumption
is highest, Hungary imports the most, $\sim$2300–2500\,MW. This pattern repeats every day and I also see it in the
forecast for week 47.

\begin{figure}[H]
\centering
\includegraphics[width=0.7\linewidth]{a4b_daily_profile_oct_nov.pdf}
\caption{Net position during the day in October and November, by year (15-min values, local time). At midday
solar almost eliminates imports; in the evening imports are highest.}
\label{fig:profil}
\end{figure}

% =====================================================================
\subsection*{b) Forecast of the net position for week 47}

\subsection*{What will drive the net position}

The forecast is based on the formula from the task: \emph{net = generation \m\ consumption} (the formula model,
labelled 10b in Figure~\ref{fig:modeli}). I therefore first looked at which parts of the formula change the most
from week to week, because for a forecast it does not matter who produces the most, but who changes the most.
Consumption, gas-fired plants, solar and Paks change the most, while the other plants (wind, lignite, biomass,
hydro) hardly change at all. For week 47 this means:
\begin{itemize}\setlength\itemsep{1pt}
\item \textbf{temperature → consumption:} the forecast average temperature for week 47 is 7.2\,°C (week 46:
6.5\,°C); the beginning of the week is warmer (up to 10.5\,°C on Wednesday), the end of the week colder
(3.5–4.7\,°C);
\item \textbf{solar:} more sun is forecast than in the cloudy week 46 (irradiation $\sim$1.3\,kWh/m$^2$ per day
against $\sim$0.9), but less than in week 45;
\item \textbf{Paks:} one reactor is in scheduled maintenance until the end of December (official notice on
ENTSO-E), so Paks will produce the same as in week 46, $\sim$1460\,MW;
\item \textbf{gas:} the most uncertain factor, because gas-fired plants only run when electricity is expensive.
When electricity in the region is cheap, they are switched off and Hungary imports more electricity from its
neighbours.
\end{itemize}

\subsection*{Approach}

To choose the best possible forecasting method, I compared 18 different models. I tested each model on past
weeks. For every week from 1 Jan 2023 to 15 Nov 2026 (202 weeks in total) it had to forecast the net position as
if that week were still unknown, i.e. using only data up to the Sunday before it. I then compared its forecast
with the actual net position. As a benchmark I used the simplest possible forecast: that the net position will be
the same as in the previous week. A good model has to be better than that. In Figure~\ref{fig:modeli} this simple
forecast has the value 1; a model with a value below 1 forecasts better, above 1 worse. I divided the models into
four groups:
\begin{itemize}\setlength\itemsep{1pt}
\item \textbf{simple}: same as last week, same as last year, mean of the last 2 or 4 weeks;
\item \textbf{one driver}: the last two weeks with a correction only for temperature, only for solar, only for
Paks or only for gas;
\item \textbf{all drivers}: linear regression with all drivers, the formula model, similar days from the past,
two neural networks;
\item \textbf{combinations}: the average of several models.
\end{itemize}

\begin{figure}[H]
\centering
\includegraphics[width=0.7\linewidth]{b3_model_comparison.pdf}
\caption{Error of the weekly forecast relative to the naive model (lower is better), test on 202 weeks
2023–2026.}
\label{fig:modeli}
\end{figure}

Test results: the best were the linear regression and the combination of the regression, the formula model and
similar days – both 34\,\% better than the naive forecast (the forecast that the week will be the same as the
previous one). The formula model from the task is right behind them (31\,\%). Of the individual drivers, solar
helps the most, while gas alone does not help at all because it is too unpredictable. The neural networks are not
better than the regression, because I have too little data for complex models ($\sim$1300 days), so they can also
learn noise. As the final model I chose the combination, because it was the best in the autumn weeks from October
to December.

The combination of models seemed the most reliable to me. It is the average of three models that reach their
forecast in different ways:
\begin{itemize}\setlength\itemsep{1pt}
\item \textbf{formula model from the task (10b):} forecasts consumption separately (based on the forecast
temperature and the day of the week) and each type of power plant, and then calculates net = generation \m\
consumption;
\item \textbf{linear regression (9):} one equation that learns from previous years how much each driver changes
the net position (temperature, solar, Paks, gas in recent days, day of the week, net position of the previous
week);
\item \textbf{similar days (11):} finds the 10 most similar days in the last two years (similar temperature and
sun, same type of day, same season) and takes the average of their net position.
\end{itemize}


\subsection*{Forecast}

For week 47 I forecast an average net position of \textbf{\m1873\,MW}, which means that Hungary will import on
average 1873\,MW at all times, $\sim$315\,GWh over the whole week. The range is from \m2146 to \m1483\,MW: it is
calculated from the model's errors in the autumn weeks of the test; in 8 out of 10 weeks the actual net position
was within this band.

Table~\ref{tab:napoved} shows the forecast by day. Average power (MW) shows how much Hungary imports on average
at every moment of the day, and energy (GWh) how much it imports in total during that day (average power × 24
hours). For the week, power is the average of the seven days and energy is their sum.

\begin{table}[H]
\centering\footnotesize
\begin{tabular}{lrrll}
\toprule
day & average power (MW) & energy (GWh) & highest import & lowest import \\
\midrule
Mon 16 Nov & \m2201 & \m52.8 & \m3121 at 19:15 & \m925 at 11:15 \\
Tue 17 Nov & \m1942 & \m46.6 & \m2862 at 19:15 & \m666 at 11:15 \\
Wed 18 Nov & \m1732 & \m41.6 & \m2652 at 19:15 & \m457 at 11:15 \\
Thu 19 Nov & \m1811 & \m43.5 & \m2731 at 19:15 & \m536 at 11:15 \\
Fri 20 Nov & \m1890 & \m45.4 & \m2810 at 19:15 & \m614 at 11:15 \\
Sat 21 Nov & \m1790 & \m43.0 & \m2710 at 19:15 & \m515 at 11:15 \\
Sun 22 Nov & \m1748 & \m41.9 & \m2668 at 19:15 & \m472 at 11:15 \\
\bottomrule
\end{tabular}
\caption{Forecast of the net position for week 47 by day (\m\ = import; local time CET).}
\label{tab:napoved}
\end{table}

\begin{figure}[H]
\centering
\includegraphics[width=0.7\linewidth]{b1_forecast_week47.pdf}
\caption{Forecast for week 47 in 15-minute steps with a 10–90\,\% range; the dotted line is the weekly average
(\m1873\,MW).}
\label{fig:napoved}
\end{figure}

I expect the highest imports on Monday, because it will be colder and cloudy, and the lowest on Wednesday, when
it will be warmest with more sun. At the weekend consumption is lower, but it will be colder, so imports will not
fall much. Every day imports are lowest around 11:00, when solar plants produce, and highest around 19:00, when
there is no sun and consumption is highest (Figure~\ref{fig:napoved}).

% =====================================================================
\subsection*{c) Will the price in week 47 be higher or lower than 290.98\,€/MWh?}

The electricity price on the exchange is set by the most expensive power plant that is still needed to cover
consumption. In Hungary this is most often a gas-fired plant, so the electricity price follows the gas price.
Because Hungary is strongly connected with its neighbours, they also shape the price: when electricity is cheap
in neighbouring countries, Hungary imports instead and the price falls.

\subsection*{Drivers from week 45 to week 47}

\begin{table}[H]
\centering\footnotesize
\footnotesize\setlength{\tabcolsep}{4pt}
\begin{tabular}{lrrrl}
\toprule
driver & week 45 & week 46 & week 47 & effect on price, w47 vs w45 \\
\midrule
wind in northern Germany (km/h) & 17.4 & 23.2 & \textbf{29.5} & \textbf{strongly down} (cheap region) \\
temperature HU (°C) & 5.3 & 6.5 & 7.2 & down: less heating \\
TTF gas price (€/MWh) & 90.2 & 87.4 & 91.0 & neutral \\
Paks (MW) & 1422 & 1458 & 1460 & neutral \\
solar HU (MW) & 653 & 288 & 426 & up: less sun than in w45 \\
residual load HU (MW) & 3349 & 3943 & 3581 & slightly up \\
price HU (€/MWh) & 291.0 & 285.9 & ? & w46 already below w45 \\
prices AT / SK / DE (€/MWh) & 287 / 293 / 257 & 280 / 286 / 245 & ? & lower everywhere in w46 \\
\bottomrule
\end{tabular}
\caption{Price drivers (weekly averages; week 47 = forecast from 14 Nov or the latest known value).}
\end{table}

In Hungary the drivers are mixed: warmer weather lowers the price, while less sun than in week 45 raises it. The
gas price is practically unchanged, so it does not affect whether the price will be higher or lower. What matters
is the wind in Germany, which will be almost twice as strong as in week 45. When the wind blows strongly in
Germany, electricity there is cheap, and because Hungary is strongly connected with its neighbours, its price
falls too. Prices in neighbouring countries already fell in week 46.

\subsection*{Check with price models}

I tested four price models on 150 weeks from 2024 onwards: how often they correctly identify whether a week will
be more or less expensive than the week two weeks earlier (as week 47 against week 45). Three of the four,
including the best in the test, forecast a lower price; the model using only Hungarian data is the worst,
because the Hungarian price is strongly determined by its neighbours.

\begin{table}[H]
\centering\footnotesize
\footnotesize\setlength{\tabcolsep}{4pt}
\begin{tabular}{lrrr}
\toprule
model & weekly error (€/MWh) & correct direction & week 47 (€/MWh) \\
\midrule
linear regression (TTF, residual load, wind DE, day) & \textbf{15.6} & 74\,\% & \textbf{263} \\
neural network (same inputs) & 16.8 & 75\,\% & 263 \\
naive (same as last week) & 17.3 & 74\,\% & 286 \\
merit order, HU data only (TTF × factor) & 20.8 & 68\,\% & 293 \\
\bottomrule
\end{tabular}
\caption{Price models: test on 150 weeks and forecast for week 47.}
\end{table}
\vspace{-8pt}

\textbf{Decision: the average price in week 47 will be lower than 290.98\,€/MWh, at about 260–275\,€/MWh.} The
same cause, cheap electricity in the region due to wind, also means that Hungarian gas-fired plants will run less
and imports will be higher.

\textbf{Risks:} severe cold, no wind, a jump in the gas price or an unplanned power plant outage. If the wind in
Germany were weaker than in week 45, the price would probably be higher.

\end{document}
